\Article{4}
Ich gehe jetzt zu \emph{bestimmten} Anordnungen der Viertelsysteme über und setze den mit $m$ bezeichneten Modul $=ea+ebi$. Dabei bezeichnet die positive ganze Zahl $e$ den größten reellen Divisor, also den größten gemeinsamen Divisor der beiden Bestandteile von $m$; somit sind $a,b$ teilerfremd. Das am einfachsten angeordnete Viertelsystem ist dasjenige, dessen Glieder $x+iy$ so beschaffen sind, dass $ax+by$ positiv und kleiner als $\frac12e(aa+bb)$ ist, während $ay-bx$ nichtnegativ und ebenfalls kleiner als $\frac12e(aa+bb)$ ist. Die letzte Bedingung ist absichtlich so ausgedrückt, dass auch die Fälle $ay-bx=0$ eingeschlossen sind. Man sieht leicht, dass es insgesamt $\frac12(e-1)$ solche Fälle gibt, nämlich

\[
\begin{array}{rlrl}
x&=a,&y&=b\\
x&=2a,&y&=2b\\
x&=3a,&y&=3b\\
\multicolumn{4}{c}{\text{usw. bis}}\\
x&=\frac12(e-1)a,&y&=\frac12(e-1)b
\end{array}
\]
also gar keine, wenn die Bestandteile von $m$ keinen gemeinsamen Divisor
\SourceLeaf{382}
haben. Nennen wir dieses Viertelsystem $k$ und diejenigen Systeme, die durch Multiplikation der zu $k$ gehörenden Zahlen mit $i,-1,-i$ entstehen, $k',k'',k'''$; oder man kann auch setzen
\[
k'=m+ik,\quad k''=(1+i)m-k,\quad k'''=im-ik
\]

So erhält man die folgende Regel, um zu entscheiden, ob eine beliebige vorgegebene, durch $m$ nicht teilbare ganze Zahl $x+iy$ einem Glied von $k,k',k''$ oder $k'''$ kongruent ist. Man kann nämlich $2(ax+by)$ in die Form $Pe(aa+bb)+Q$ und $2(ay-bx)$ in die Form $Re(aa+bb)+S$ bringen, wobei $P,Q,R,S$ ganze reelle Zahlen sind und folgende Bedingungen erfüllen:

\begin{center}wenn $x+iy$ einer Zahl aus folgendem System kongruent ist\end{center}
\begingroup\small\setlength{\arraycolsep}{3pt}
\[
\begin{array}{l||l|l|l|l}
\text{so dass}&k&k'&k''&k'''\\
P&\text{gerade}&\text{ungerade}&\text{ungerade}&\text{gerade}\\
R&\text{gerade}&\text{gerade}&\text{ungerade}&\text{ungerade}\\
Q&\text{positiv}&\text{positiv}&\text{positiv}&\text{nichtnegativ}\\
e(aa+bb)-Q&\text{positiv}&\text{nichtnegativ}&\text{positiv}&\text{positiv}\\
S&\text{nichtnegativ}&\text{positiv}&\text{positiv}&\text{positiv}\\
e(aa+bb)-S&\text{positiv}&\text{positiv}&\text{nichtnegativ}&\text{positiv}
\end{array}
\]
\endgroup

Die Übersicht wird leichter, wenn man die Fälle, in denen keine der Zahlen $ax+by$, $ay-bx$ durch $e(aa+bb)$ teilbar ist, von den übrigen unterscheidet.

I. Im ersten Fall nimmt man für $P$ einfach die (algebraisch) kleinere der beiden ganzen Zahlen, zwischen denen $\frac{2(ax+by)}{e(aa+bb)}$ ausschließlich liegt, und ebenso für $R$ die kleinere der beiden, zwischen denen $\frac{2(ay-bx)}{e(aa+bb)}$ liegt.

II. Ist $ax+by$ durch $e(aa+bb)$ teilbar, so ist $ay-bx$ zwar durch $aa+bb$, aber nicht durch $e(aa+bb)$ teilbar (weil sonst $x+iy$ durch $ea+ebi$ teilbar wäre). Ist nun $R$, also die unmittelbar unter $\frac{2(ay-bx)}{e(aa+bb)}$ liegende ganze Zahl, gerade, so ist $x+iy$ modulo $ea+ebi$ einer Zahl aus $k'$ kongruent, dagegen einer aus $k'''$, wenn $R$ ungerade ist.

III. Ist $ay+bx$ durch $e(aa+bb)$ teilbar, nicht aber $ax+by$, so ist $x+iy$ einer Zahl aus $k$ oder aus $k''$ kongruent, je nachdem die algebraisch unmittelbar unter $\frac{2(ax+by)}{e(aa+bb)}$ liegende ganze Zahl gerade oder ungerade ist.

\SourceLeaf{383}

\Article{5}
Aus dem Vorstehenden ergibt sich ohne Mühe die folgende Methode, um den Charakter einer gegebenen ganzen Zahl $M=A+Bi$ bezüglich des ungeraden Moduls $m=ea+ebi$, der sie nicht teilt, zu bestimmen.

Zur Abkürzung verwenden wir folgende Bezeichnung. Ist $p$ eine beliebige nichtganze reelle Zahl, so bezeichnet $[p]$ diejenige ganze Zahl, die sowohl $p-[p]$ als auch $1+[p]-p$ positiv macht. Diese Definition schließt also die Anwendung der Bezeichnung auf ganze Zahlen aus. Würde man verlangen, dass weder $p-[p]$ noch $1+[p]-p$ negativ ist, so wäre das Zeichen $[p]$ für ganzzahliges $p$ zweideutig. Man könnte auch, wie in einer früheren Abhandlung geschehen, verlangen, dass $1+[p]-p$ positiv und $p-[p]$ lediglich nichtnegativ ist. Für unseren Zweck ist es etwas bequemer, bei der ersten Definition zu bleiben.

Das Viertelsystem $k$ besteht demnach aus allen ganzen Zahlen $f$, für die, wenn man $\frac{2f}{m}=\xi+\eta i$ setzt, $\xi$ ausschließlich zwischen $0$ und $1$ liegt und $\eta$ zwischen $0$ und $1$, mit eingeschlossener $0$; oder
\[
\begin{array}{lll}
&[\xi]=0,&[\eta]=0\\
\text{oder}&[\xi]=0,&\eta=0
\end{array}
\]

Setzt man nun für jedes $f$ $\frac{2fM}{m}=\xi+i\eta$ und nimmt
\begingroup\small\setlength{\arraycolsep}{2pt}
\[
\begin{array}{rl}
\Psi f=0&\text{wenn zugleich }[\xi]\text{ gerade und entweder }[\eta]\text{ gerade oder }\eta\text{ ganz}\\
1&[\eta]\text{ gerade und entweder }[\xi]\text{ ungerade oder }\xi\text{ ganz}\\
2&[\xi]\text{ ungerade und entweder }[\eta]\text{ ungerade oder }\eta\text{ ganz}\\
3&[\eta]\text{ ungerade und entweder }[\xi]\text{ gerade oder }\xi\text{ ganz}
\end{array}
\]
\endgroup

oder in Tabellenform
\[
\begin{array}{l|c|c|c}
&[\eta]\text{ gerade}&[\eta]\text{ ungerade}&\eta\text{ ganz}\\\hline
[\xi]\text{ gerade}&0&3&0\\
\relax[\xi]\text{ ungerade}&1&2&2\\
\xi\text{ ganz}&1&3&-\ \pdot
\end{array}
\]
was man mit $\nabla(\xi+i\eta)$ bezeichnen kann, so ist der gesuchte Charakter der Zahl $M$ bezüglich des Moduls $m$ gleich dem kleinsten Rest von $\Sigma\Psi f$ modulo $4$.

\SourceLeaf{384}

\Article{6}
Die im vorhergehenden Artikel gegebene Vorschrift ist allgemein; für den Fall, dass $M$ ungerade ist, werden wir ihr jedoch eine andere Gestalt geben. Zugleich nehmen wir an, dass die reellen Teile von $m$ und $M$ ungerade und daher die imaginären Teile gerade sind.

Wir ordnen jeder zum Viertelsystem $k$ gehörenden Zahl $f$ eine andere Zahl $g$ zu, die aus $f$ folgendermaßen abgeleitet wird. Setzt man $\frac{2f}{m}=\xi+i\eta$, so unterscheidet man vier Fälle:
\[
\begin{array}{ll}
\mathrm{I.}&\text{Wenn }[2\xi]=0\text{ und entweder }[2\eta]=0\text{ oder }\eta=0\\
\mathrm{II.}&\text{Wenn }[2\xi]=1\text{ und entweder }[2\eta]=0\text{ oder }\eta=0\\
\mathrm{III.}&\text{Wenn }[2\xi]=0\text{ und }[2\eta]=1\\
\mathrm{IV.}&\text{Wenn }[2\xi]=1\text{ und }[2\eta]=1
\end{array}
\]
Im ersten Fall setzt man $g=2f$, im zweiten $g=im-2if$, im dritten $g=m+2if$, im vierten $g=(1+i)m-2f$. Man sieht leicht, dass die Gesamtheit aller $g$ ein vollständiges Viertelsystem $l$ bildet. Seine kennzeichnende Eigenschaft ist, dass, wenn man $\frac gm=\xi+i\eta$ setzt, zugleich folgende Bedingungen gelten:

\[
\begin{array}{llll}
\text{entweder}&[\xi]=0,&[\eta]=0&\\
\text{oder}&\eta=0,&[\xi]=0&\text{und }g\text{ durch }1+i\text{ teilbar}\\
\text{oder}&\xi=0,&[\eta]=0&\text{und }g\text{ durch }1+i\text{ nicht teilbar}
\end{array}
\]
Daraus folgt, dass $l$ sich von $k$ nur dadurch unterscheidet, dass diejenigen Zahlen in $k$, für die $\eta=0$ und die durch $1+i$ nicht teilbar sind, nämlich
\[
\begin{aligned}
&a+bi,\quad3(a+bi),\quad5(a+bi)\ldots\frac{e-3}{2}(a+bi)\\
&\text{oder bis }\frac{e-1}{2}(a+bi)
\end{aligned}
\]
je nachdem $e$ die Form $4n+1$ oder $4n+3$ hat, in $l$ fehlen und stattdessen in der letzteren Gesamtheit die Produkte dieser Zahlen mit $i$ auftreten. Zugleich sieht man, dass für $e=1$, also wenn $m$ durch keine reelle ganze Zahl teilbar ist, $k$ und $l$ völlig gleich sind.

Nun ist $\Psi f$ unmittelbar aus dem zu $f$ gehörenden $g$ abzuleiten. Das Ergebnis ist, dass für die obigen $4$ Fälle gilt:
\[
\begin{aligned}
\mathrm{I.}\quad&\Psi f=\nabla\frac{gM}{m}\\
\mathrm{II.}\quad&\Psi f=\left\{\begin{array}{ll}
-\nabla\frac{gM}{m}&\begin{gathered}\text{wenn weder reeller noch imaginärer Teil}\\\text{ von }\frac{gM}{m}\text{ ganz}\end{gathered}\\
1-\nabla\frac{gM}{m}&\text{wenn einer von beiden ganz ist}
\end{array}\right.
\end{aligned}
\]
III. wie II.\quad IV. wie I.
\BookRule{.3\textwidth}
